
\section{NetAlly: 네트워크 운영 Multi-Agent 시스템}
\subsection{설계 동기: 왜 도메인 학습이 아닌 3-Plane 통합인가}
L4/L5의 성능 격차를 해소하는 접근으로 도메인 특화 학습(fine-tuning)과 기존 분석 도구 통합이 있다. 우리는 후자를 채택하였다. 네트워크 설정 파일은 벤더와 OS 버전에 따라 구문이 상이하여 학습 데이터 확보가 어렵고, L4/L5 태스크는 그래프 알고리즘의 실행을 요구하므로 학습량을 늘려도 정확도 상한이 존재한다. 반면 Batfish와 같이 이미 검증된 분석 도구를 활용하면 정확도와 신뢰성을 동시에 확보할 수 있다.

\begin{figure}[!t]
  \centering
  \includegraphics[width=0.5\textwidth]{figure/netally_arch.png}
  \caption{NetAlly 3-Plane Architecture. The Orchestrator analyzes the query to select relevant skills, and the Executor invokes tools across the PNETLab, NSO, and Batfish planes through a ReAct loop.}
  \label{fig:netally-arch}
\end{figure}


\subsection{3-Plane 통합 아키텍처}
Fig.~\ref{fig:netally-arch}에 나타낸 바와 같이, NetAlly는 세 개의 독립된 네트워크 관리 평면을 하나의 Chat 인터페이스로 통합한다: (1) \textbf{PNETLab Plane} --- 가상 토폴로지 관리, 장비 부팅, 콘솔 접근, (2) \textbf{NSO Plane} --- RESTCONF 기반 구성 조회\textperiodcentered 동기화\textperiodcentered 변경, (3) \textbf{Batfish Plane} --- 정형 검증, traceroute 시뮬레이션, What-If 분석. 이 세 평면의 도구들은 MCP를 통해 에이전트에 노출되며, Orchestrator--Executor 구조가 LangGraph~\cite{ref3} 기반 상태 기계(State Machine) 위에서 질의의 성격에 따라 적절한 평면의 도구를 선택한다.

\textbf{Orchestrator (계획 에이전트):} 경량 LLM이 스킬 카탈로그(각 스킬의 이름, 설명, 필요 도구가 YAML 프론트매터로 정의된 SKILL.md 파일 집합)와 사용자 질의를 입력받아, JSON 형식으로 1-3개의 관련 스킬을 선택하고 선택 근거를 출력한다. 예를 들어, ``P1-P2 링크 장애 시 PE1$\rightarrow$Leaf3 경로는?''이라는 질의에 대해 Orchestrator는 \texttt{["core", "network\_verify"]}를 선택하고, ``What-If 분석이 필요하므로 Batfish fork\_snapshot을 사용해야 한다''는 추론을 생성한다. 선택된 스킬의 상세 가이드(SKILL.md 본문)가 Executor의 시스템 프롬프트에 주입되어, 도구 호출의 맥락과 제약 조건을 전달한다.

\textbf{Executor (실행 에이전트):} Executor LLM은 Orchestrator가 주입한 스킬 프롬프트와 16개 이상의 MCP 도구 바인딩을 기반으로 ReAct(Reason-Act-Observe) 루프를 수행한다. 각 반복에서 대화 이력을 분석하여 다음 행동을 결정하고, 도구를 호출한 뒤 출력을 관찰하여 추가 호출 여부를 판단한다. 최종 답변이 도출되거나 최대 10단계에 도달하면 루프가 종료된다.


\subsection{다층 오류 복구 메커니즘}
NetAlly는 다층 오류 복구(Multi-layer Error Recovery)를 통해 단순 도구 래퍼(Wrapper)와 차별화된다. 도구 호출 실패를 Python 예외로 전파하지 않고, 구조화된 오류 딕셔너리(\texttt{\{"error": ..., "tool": ...\}})로 변환하여 LLM의 대화 이력에 삽입한다. 이를 통해 LLM은 다음 반복에서 (a)~다른 도구로 전환, (b)~동일 도구를 수정된 파라미터로 재호출, (c)~부분적 결과로 최선의 답변 생성 중 하나를 선택한다. 이 외에도 \texttt{async} 호출 실패 시 \texttt{sync} 스레드로의 자동 폴백, \texttt{NETALLY\_MCP\_ALLOW\_MUTATIONS} 플래그를 통한 읽기 전용 모드 강제, 10회 반복 후 강제 종료를 통한 무한 루프 방지, 미포착 예외의 SSE \texttt{error} 이벤트 변환 등의 안전 장치를 포함한다. 예를 들어, Batfish가 특정 설정 파일을 파싱하지 못하는 경우에도 에이전트는 NSO 쿼리로 대안을 시도하며, 완전한 실패보다 부분적 답변을 우선한다.

\subsection{핵심 도구와 자율 파이프라인}
Executor가 호출하는 핵심 도구는 세 가지이다: (1)~\texttt{network\_query}는 Cisco NSO RESTCONF를 통해 실시간 구성을 조회하며 L1- L3에 주로 활용된다. (2)~\texttt{network\_verify}는 Batfish를 통해 정형 검증과 What-If 분석을 수행하며 L4- L5에 활용된다. (3)~\texttt{lab\_manage}는 PNETLab 토폴로지 관리 및 자동 온보딩을 담당한다.
NetAlly는 두 가지 자동화 파이프라인을 통해 수동 조작을 최소화한다.

\textbf{자율 온보딩.} Section~V의 실험에서 4개 Lab의 설정 파일 수집과 Batfish 스냅샷 초기화에 사용된다. PNETLab의 가상 장비를 사람의 CLI 조작 없이 운영 가능 상태로 전환한다. 파이프라인은 (1)~LabFS에서 토폴로지와 콘솔 포트를 파싱하고, (2)~각 장비에 Telnet으로 접속하여 SSH를 활성화하며, (3)~NSO에 RESTCONF로 등록 후 설정을 동기화하고, (4)~수집된 설정으로 Batfish 스냅샷을 초기화한다. IOS 버전별 명령어 호환성, 관리 IP 충돌, 동기화 실패 시 자동 재시도 등의 예외 처리가 포함된다.

\textbf{Fork-Snapshot 기반 What-If 분석.} L5 장애 영향 추론을 지원하기 위해, Batfish의 \texttt{fork\_snapshot}으로 지정된 변경(링크/노드 비활성화)을 적용한 가상 네트워크 복사본을 생성한다. 원본 스냅샷은 보존되므로 비파괴적이다. 생성된 복사본에 \texttt{differentialReachability} 쿼리를 실행하여, 장애 전후의 도달성 차이를 식별하고 영향 범위를 도출한다. 이 과정은 에이전트가 수행하며, 운영자는 ``P1-P2 링크가 다운되면 어떤 영향이 있는가?''와 같은 자연어 질의만 제공하면 된다.